
\textbf{Proposition 3.4.} \uline{Avec les hypothèses et les notations de} 3.1 \uline{et} 3.3, \uline{on a pour tout} \(x\in K^{\bullet}(Y)\) \uline{la relation}:
\[
f^*(i_*(x))=j_*\bigl(g^*(x).\lambda_{-1}(F)\bigr).
\]

Comme les deux membres de la relation à démontrer sont additifs, on peut supposer que \(x\) est la classe d'un \(\underline{O}_Y\)-Module localement libre de type fini \(\underline{M}\). L'élément \(f^*(i_*(x))\) est alors la classe du complexe \(\underline{\underline{L}}f^*(\underline{\underline{R}}i_*(\underline{M}))\), obtenu en prenant une résolution projective de \(i_*(\underline{M})\) et en lui appliquant le foncteur \(f_*\). Les faisceaux d'homologie du complexe ainsi obtenu sont les \(\underline{\Tor}_i^{\underline{O}_X}(\underline{M},\underline{O}_{X'})\). Or d'après 2.4 et 3.2, on a:
\[
\underline{\Tor}_i^{\underline{O}_X}(\underline{M},\underline{O}_{X'})
\xrightarrow{\sim}
\underline{M}\otimes_{\underline{O}_Y}\underline{\Tor}_i^{\underline{O}_X}(\underline{O}_{X'},\underline{O}_Y)
\xrightarrow{\sim}
\underline{M}\otimes_{\underline{O}_Y}\bigwedge^i\underline{F}.
\]

Il en résulte que ces faisceaux d'homologie sont localement libres de type fini sur \(\underline{O}_{Y'}\), donc parfaits sur \(\underline{O}_{X'}\); par suite, la classe dans \(K^{\bullet}(X')\) du complexe \(\underline{\underline{L}}f^*(\underline{\underline{R}}i_*(\underline{M}))\) est la somme alternée des classes de ses faisceaux d'homologie, soit:
\[
f^*(i_*(x))
=\sum_{i\ge0}(-1)^i\operatorname{cl}
\bigl(j_*(\underline{M}\otimes_{\underline{O}_Y}\bigwedge^i\underline{F})\bigr)
=\sum_{i\ge0}(-1)^i
j_*\bigl(g^*(x)\lambda^i(F)\bigr).
\]

\textbf{Lemme 3.5.} \uline{Avec les hypothèses et les notations de} 3.1, \uline{on a les isomorphismes suivants} :
\begin{enumerate}[label=\roman*)]
\item \uline{pour tout} \(n\ge0\), \uline{et pour tout} \(i\ge1\), \(R^if_*\bigl(\underline{O}_{X'}(n)\bigr)=0\);
\item \uline{pour tout} \(n\ge0\), \(f_*\bigl(\underline{O}_{X'}(n)\bigr)=\underline{J}^n\).
\end{enumerate}

On peut évidemment supposer \(X\) affine, donc quasi-compact.

i) Rappelons d'abord le lemme suivant, qui a été démontré dans exp.

